\chapter{Glossary}
\label{ch:glossary}

\begin{description}

\item[catalog] The imported report register --- report code, vehicle, title,
discipline, date, authors --- independent of whether the file behind a row has
been ingested.

\item[catalog link] The join between a catalog row and an ingested document. At
most one document per catalog entry.

\item[chunk] The unit of retrieval: a block of about 650 words with 75 words
overlapping its neighbour, carrying the page range and section title it came
from. Search, review and comparison all operate on chunks, never on raw pages.

\item[discipline profile] A set of review rules for one engineering discipline
(NVH, CFD, Durability, Test/Validasyon), defined as data in
\file{app/rules/profiles/<name>.json}.

\item[document] One ingested file, identified by the SHA-256 hash of its bytes.

\item[extraction quality] The per-document JSON rollup describing how well text
came out: pages parsed and stored, OCR counts, sparse and unreadable pages, mean
characters per page.

\item[finding] One complaint from the review engine, carrying a rule id, a
severity, the evidence it rests on, the page range, and a suggested fix.

\item[finding key] A hash over the identifying parts of a finding. A human
decision is recorded against this key, so the decision survives a re-run.

\item[hybrid search] The default retrieval mode: keyword and semantic results
normalized and combined at 0.45 / 0.55, with a lexical rerank contribution and a
guard against semantic-only matches that share no query token.

\item[job] A background unit of work with an id, a status, progress and a
result. Jobs run on a single worker thread; poll
\route{GET /jobs/\{job\_id\}}.

\item[near-duplicate] Two documents that are the same report in substance but not
byte identical. Detected by title and embedding similarity, unlike ingestion's
exact hash match.

\item[provider] The named implementation behind a pluggable role --- an embedding
provider (\code{token-hash-v1}, \code{sentence-transformers:<model>}), a
retrieval provider, a generation provider. Reported in responses so an answer
can be traced to what produced it.

\item[retrieval version] Which retrieval implementation answered: \code{v1}
(legacy symmetric embedding with a strict token gate), \code{v2} (the default
asymmetric in-house path) or \code{v3} (the Haystack pipeline). A request
parameter, not configuration.

\item[section] What a parser produces per unit of a document: a PDF page, a DOCX
heading-delimited block, or a PPTX slide. Sections become
\code{document\_pages} rows after cleaning.

\item[selective OCR] OCR applied only to PDF pages whose selectable text falls
below \env{OCR\_MIN\_TEXT\_CHARACTERS}, merging the OCR text with whatever native
text the page had.

\item[skill] A packaged capability outside the ordinary request path. The CATIA
mass/CG skill ships as \file{skill/catia-mass-cg.skill} and carries its own
harness and safety gates.

\item[token-hash] The deterministic fallback embedding provider: a
256-dimension hashed bag-of-tokens vector needing no model, no GPU and no
network. Lower retrieval quality, identical interface --- which is what the test
suite depends on.

\item[variant] The product identity selected by \env{APP\_VARIANT}:
\code{big\_agent} (SmartCAE AI), \code{raporhub} or \code{repocto}. It changes
branding and the served front end, not the engines.

\item[vector index] The process-wide NumPy matrix of all stored embeddings,
stamped with a cheap database signature so it is rebuilt whenever embeddings
change.

\end{description}
